\section{Proof of auxiliary nonvanishing}
\label{sec:auxiliary-nonvanishing}

We prove Proposition~\ref{prop:rough-nonvanishing}. The even-sign
assertion follows from Hoffstein--Luo after one preliminary twist.
For the odd-sign assertion we retain the prescribed local conditions
throughout a derivative first-moment calculation.

\begin{proof}[Proof of Proposition~\ref{prop:rough-nonvanishing}\textup{(i)}]
An odd fundamental discriminant is squarefree and congruent to $1$
modulo $4$, so $c_2$ is represented by $u_2\in\{1,5\}$ modulo $8$.
For odd $p\in S$, choose a nonzero residue $u_p$ representing $c_p$.
The Chinese remainder theorem and Dirichlet's theorem give a prime
$q\notin S\cup B$ satisfying
\[
 q\equiv\sigma u_2\pmod8,\qquad
 \sigma q\equiv u_p\pmod p\quad(p\in S,\ p\ne2).
\]
Then $D_0=\sigma q$ is an odd fundamental discriminant in the prescribed
filter. Put $A=E^{(D_0)}$ and let $N_A$ be its conductor. Thus
$\varepsilon(A)=1$.

The theorem of Hoffstein--Luo
\citep[pp.~435--436]{HoffsteinLuo1997} gives infinitely many odd
fundamental discriminants $h$ such that
\[
 L(A^{(h)},1)\ne0,\qquad \omega(|h|)\le4,
\]
and $\chi_h(p)=1$ at every prime in the fixed finite set consisting of
$S\cup B\cup\{q\}$ and the prime divisors of $N_A$.
The twisting formula for the root number gives
\[
 1=\varepsilon(A^{(h)})=\varepsilon(A)\chi_h(-N_A)
 =\operatorname{sign}(h),
\]
so $h>0$. For odd $p$ the condition $\chi_h(p)=1$ says that $h$ is a
unit square in $\Q_p$. At $2$, odd fundamentality and $\chi_h(2)=1$
give $h\equiv1\pmod8$, which has the same meaning.
Since $q\nmid h$, the product $D=D_0h$ is squarefree and congruent to
$1$ modulo $4$. It retains the prescribed sign and squareclasses and
avoids $B$. Finally,
\[
 E^{(D)}\simeq A^{(h)},\qquad \omega(|D|)=1+\omega(h)\le5.
\]
Distinct choices of $h$ give distinct $D$.
\end{proof}

\subsection{An odd-sign witness in a fixed progression}
\label{odd:fixed-progression}

We retain the modulus of the local conditions throughout the first-moment
calculation.  Iwaniec's derivative moment
\citep[Theorem on p.~367 and Sections~6--10]{Iwaniec1990} treats the
negative-discriminant split family.  The following argument allows any
fixed admissible progression.  It makes no uniformity assertion as that
progression varies.

\begin{lemma}[Odd first moment]\label{nv:odd-first-moment}
Let $f$ be the weight-two newform of an elliptic curve of conductor $N$,
with its $L$-function normalized to have center $1/2$.
Fix $H$ divisible by $8$ and $N$, a reduced class $a\pmod H$, and
$\sigma\in\{1,-1\}$ with $\sigma a\equiv1\pmod4$.
Suppose the twists with $n>0$ squarefree and $n\equiv a\pmod H$
have root number $-1$.
For every nonnegative, nonzero $\Phi\in C_c^\infty((1,2))$, there are
$c>0$, $c_1\in\mathbb R$, and an absolute $\eta>0$ such that
\begin{equation}\label{odd:moment}
 \sum_{\substack{n>0\ \mathrm{squarefree}\\n\equiv a\ (H)}}
 L'(1/2,f\otimes\chi_{\sigma n})\Phi(n/X)
 =cX\log X+c_1X+O_{f,H,a,\sigma,\Phi}(X^{1-\eta}).
\end{equation}
Consequently this progression contains arbitrarily large $n$ whose
twists have a simple central zero.
\end{lemma}

\begin{proof}
Write $L(s,f)=\sum_{m\ge1}\lambda(m)m^{-s}$, so
$|\lambda(m)|\le\tau(m)$, where $\tau$ is the divisor function, and put $Q=\sqrt N/(2\pi)$.
In estimates below, constants may depend on the fixed data in the Lemma;
all unspecified power exponents are absolute.
The proof separates two issues: preserving the fixed progression while
removing the squarefree condition, and controlling the frequencies after
Poisson summation. The zero frequency will give a positive main term of
order $X\log X$; the other frequencies and the discarded square divisors
must together have a power-saving bound. We keep the modulus $H$ fixed
throughout these steps.
Set
\[
 V(y)=\frac1{2\pi i}\int_{(2)}
 e^{z^2}\Gamma(1+z)y^{-z}\frac{dz}{z^2}.
\]
The conductor of the twist is $Nn^2$.  Its odd functional equation gives
\begin{equation}\label{odd:afe}
 L'(1/2,f\otimes\chi_{\sigma n})
 =2\sum_{m\ge1}\frac{\lambda(m)\chi_{\sigma n}(m)}{\sqrt m}
           V\left(\frac m{Qn}\right).
\end{equation}
Indeed, shift the defining completed integral across its double pole
and use the functional equation.  The central value is zero, so the
derivative of the completion contributes no additional term.
Use the right side of Equation~\eqref{odd:afe} to define $\mathcal A(n)$
also for nonsquarefree $n\equiv a\pmod H$.

\paragraph{Large square divisors.}
We first record the precise large-sieve input.  For every $\rho>0$ and
fixed $0<\delta<1/4$,
\begin{equation}\label{odd:second-moment}
 \sum_{\substack{0<t\le Z\ \mathrm{squarefree}\\
             \sigma t\ \mathrm{fundamental},\ (t,2N)=1}}
 |L(1/2+\delta+i v,f\otimes\chi_{\sigma t})|^2
 \ll_{f,\rho,\delta} Z^{1+\rho}(1+|v|)^{C_{\rho,\delta}}.
\end{equation}
To justify this consequence of the quadratic large sieve, its form in
\citep[Theorem~1]{HeathBrown1995} is
\[
 \sum_{\substack{t\le Z\ \mathrm{odd, squarefree}}}
 \left|\sum_{\substack{u\le U\ \mathrm{odd, squarefree}}}
 b_u\left(\frac tu\right)\right|^2
 \ll_\epsilon (ZU)^\epsilon(Z+U)\sum_u|b_u|^2.
\]
Quadratic reciprocity permits either orientation of the symbol after
separating classes modulo $4$.  In a Dirichlet polynomial with coefficients
$\lambda(m)m^{-1/2}$, write $m=ur^2$, $u$ squarefree.  For each $r$,
$\chi_{\sigma t}(r^2)$ restricts $t$ to $(t,r)=1$; discard that
restriction only after taking absolute squares.  Since
$|\lambda(ur^2)|\le\tau(u)\tau(r^2)$, Minkowski's inequality bounds the
square root of the moment by
\begin{equation}\label{odd:square-index-bound}
 (ZU)^{O(\epsilon)}
 \sum_{r\le\sqrt U}\frac{\sqrt{Z+U/r^2}}r
 \ll (ZU)^{O(\epsilon)}(\sqrt Z+\sqrt U).
\end{equation}
The powers of primes dividing $2N$ contribute convergent geometric sums.
On a dyadic interval $t\asymp T$, the approximate functional equation
has length $U\ll_f T(1+|v|)^C$, apart from rapidly decreasing tails.
The main sum's additional factor $m^{-\delta}$ improves partial
summation.  The dual sum loses at most $U^\delta$ in coefficients, but
its functional-equation multiplier contributes $T^{-2\delta}$, up to
a fixed power of $1+|v|$.  Thus it incurs no positive power of $T$.
Mellin inversion separates the smooth conductor weights; the gamma
factors give integrable decay in the Mellin variables.  Applying
Equation~\eqref{odd:square-index-bound} with sufficiently small
$\epsilon$, and summing dyadic intervals, proves
Equation~\eqref{odd:second-moment}.

Expand $\mu^2(n)=\sum_{d^2\mid n}\mu(d)$ and retain $d\le Y$, where
$1\le Y\le\sqrt X$.  The omitted contribution is
\begin{equation}\label{odd:square-tail}
 O_\rho(X^{1+\rho}/Y).
\end{equation}
Here are the details for the nonsquarefree terms.  Write $n=tv^2$ with
$t$ squarefree.  An omitted term has $v>Y$, with at most $\tau(v)$ choices
of $d$.  Also $\sigma t$ is fundamental, $(tv,2N)=1$, and
\[
 \chi_{\sigma n}(m)
   =\chi_{\sigma t}(m)\mathbf1_{(m,v)=1}.
\]
On a Mellin line $\Re z=\delta>0$, the series defining $\mathcal A(tv^2)$
therefore uses
\[
 L(1/2+z,f\otimes\chi_{\sigma t})
 \prod_{p\mid v}
 \left(1-\frac{\lambda(p)\chi_{\sigma t}(p)}{p^{1/2+z}}
             +\frac{\chi_{\sigma t}(p)^2}{p^{1+2z}}\right).
\]
The product is $O_\epsilon(v^\epsilon)$.  By Cauchy's inequality and
Equation~\eqref{odd:second-moment}, the sum of the absolute $L$-values
over $t\le2X/v^2$ is
$O((X/v^2)^{1+\epsilon}(1+|\Im z|)^C)$.
The remaining Mellin factors are bounded by $O(X^\delta)$ times a
rapidly decreasing function of $\Im z$.  Choose $\delta$ and $\epsilon$
sufficiently small in terms of $\rho$.  Absorbing the divisor count and
these small powers gives
$X^{1+\rho}\sum_{v>Y}v^{-2}\ll X^{1+\rho}/Y$, as claimed.
The double pole in the kernel is harmless on this fixed positive line.

\paragraph{Poisson summation.}
For $(d,H)=1$ put $J=d^2$ and $c_J\equiv aJ^{-1}\pmod H$.
Let $T(J)$ be the sum of $\mathcal A(n)\Phi(n/X)$ over
$n\equiv a\pmod H$ with $J\mid n$, without a squarefree restriction.
Writing $m=m_1q$ with $m_1\mid H^\infty$ and $(q,H)=1$, put
$\kappa(m_1)=\chi_{\sigma a}(m_1)$ and
\[
 G_\nu(q)=\sum_{y\bmod q}\left(\frac yq\right)e(\nu y/q),\qquad
 W(z,R)=\int_0^\infty\Phi(x)V\left(\frac z{Qx}\right)e(-Rx)\,dx,
\]
where $e(x)=e^{2\pi ix}$. Poisson summation gives
\begin{equation}\label{odd:poisson}
 \begin{split}
 T(J)={}&\frac{2X}{HJ}
 \sum_{m_1\mid H^\infty}\frac{\lambda(m_1)\kappa(m_1)}{\sqrt{m_1}}
 \sum_{(q,H)=1}\frac{\lambda(q)}{q^{3/2}}
                  \left(\frac{\sigma JH}{q}\right)\\
 &\hspace{9mm}\cdot\sum_{\nu\in\mathbb Z}
 G_\nu(q)e(\nu c_J\bar q/H)
 W\left(\frac{m_1q}{X},\frac{\nu X}{HJq}\right),
 \end{split}
\end{equation}
where $q\bar q\equiv1\pmod H$ and $G_\nu(1)=1$.
For clarity, the Chinese remainder theorem evaluates the finite sum as
\[
 \sum_{\substack{x\bmod Hq\\x\equiv c_J\ (H)}}
 \left(\frac xq\right)e(\nu x/(Hq))
 =e(\nu c_J\bar q/H)\left(\frac Hq\right)G_\nu(q).
\]
Multiplication by $(\sigma J/q)$ yields Equation~\eqref{odd:poisson};
this factor vanishes when $(q,J)>1$.

\paragraph{The zero frequency.}
The Gauss sum $G_0(q)$ is zero unless $q$ is a square, when it is
$\varphi(q)$.  Define
\[
 H_p(z)=1+(1-p^{-1})\sum_{r\ge1}
                    \lambda(p^{2r})p^{-r(1+2z)},\qquad
 F_J(z)=\prod_{p\nmid HJ}H_p(z),
\]
\[
 L_H(z)=\sum_{m_1\mid H^\infty}
          \lambda(m_1)\kappa(m_1)m_1^{-1/2-z},\qquad
 \widehat\Phi_0(z)=\int_0^\infty\Phi(x)x^z\,dx.
\]
Mellin inversion expresses the zero-frequency contribution as
\begin{equation}\label{odd:zero-integral}
 T_0(J)=\frac{2X}{HJ}\frac1{2\pi i}\int_{(2)}
 e^{z^2}\Gamma(1+z)(QX)^z\widehat\Phi_0(z)L_H(z)F_J(z)
 \frac{dz}{z^2}.
\end{equation}
At a good prime, let $\alpha_p\beta_p=1$ and
$|\alpha_p|=|\beta_p|=1$.  The identity
\[
 \sum_{r\ge0}\lambda(p^{2r})x^r
 =\frac{1+x}{(1-\alpha_p^2x)(1-\beta_p^2x)}
\]
shows that $H_p(z)$ is the symmetric-square Euler factor at $1+2z$
times a polynomial
$1+O(p^{-2-4\Re z}+p^{-2-2\Re z})$.
The correction product converges absolutely near $\Re z\ge-1/16$.
To omit a prime dividing $J$, omit its correction factor and insert an
inverse symmetric-square Euler polynomial; this introduces no poles at
zeros of $H_p$.  These omitted factors cost at most $J^B$ on that strip.
The finite Euler product $L_H$ is holomorphic there.
Symmetric-square continuation, its functional equation, and polynomial
strip bounds now allow the shift to $\Re z=-1/16$, with error
\begin{equation}\label{odd:zero-error}
 T_0(J)=R_J(X)+O(J^B X^{1-1/16}),
\end{equation}
where $R_J(X)$ is the residue of the integrand in
Equation~\eqref{odd:zero-integral}, including its factor $2X/(HJ)$.
For this analytic input see
\citep[Theorem~9.3 and Remark~9.9]{GelbartJacquet1978}; in the CM case
the untwisted symmetric square is a product of nontrivial Dirichlet and
Hecke-character $L$-functions, so the same continuation holds near $1$.

The sum of the residues, with $d$ extended to infinity, has coefficient
function
\[
 M(z)=e^{z^2}\Gamma(1+z)Q^z\widehat\Phi_0(z)L_H(z)
              \prod_{p\nmid H}\bigl(H_p(z)-p^{-2}\bigr).
\]
The product is interpreted by its symmetric-square factorization and
an absolutely convergent correction.  It is holomorphic near zero.
Moreover $M(0)>0$.  Indeed
\[
 F_p^{\rm ev}:=\frac{1+p^{-1}}
 {(1-\alpha_p^2/p)(1-\beta_p^2/p)}>0,\qquad
 H_p(0)=p^{-1}+(1-p^{-1})F_p^{\rm ev}>p^{-1}.
\]
All the fixed local factors in $L_H(0)$ are positive, and
$L(1,\operatorname{sym}^2 f)>0$ by the Rankin--Selberg residue formula.
The remaining correction factors are positive and form a convergent
product.  The full residue is therefore
\begin{equation}\label{odd:residue}
 \frac{2X}{H}\bigl(M(0)\log X+M'(0)\bigr),
 \qquad c=\frac{2M(0)}H>0.
\end{equation}
Extending the residue sum costs $O_\rho(X^{1+\rho}/Y)$:
its terms contain $d^{-2}\prod_{p\mid d}H_p(0)^{-1}$, and this
product, as well as its differentiated version, is
$O_\epsilon(d^\epsilon(1+\log d))$.
The resulting $d^{-2+\epsilon}$ tail also justifies differentiating
the residue sum.

\paragraph{The nonzero frequencies.}
The zero frequency has now supplied the positive main term. We next
bound the remaining frequencies uniformly in the square divisor $J$
by a fixed power of $J$, so that summation over the retained square
divisors will preserve a power saving in $X$.
We prove the bound
\begin{equation}\label{odd:offdiagonal}
 T_{\ne0}(J)\ll J^B X^{7/8}.
\end{equation}
For odd $q$, set $\epsilon_q=1$ or $i$ according as $q\equiv1$ or
$3\pmod4$, and $\widetilde G_\nu(q)=\epsilon_q^{-1}G_\nu(q)$.
Quadratic reciprocity makes $\widetilde G_\nu$ multiplicative in $q$.
As $4\mid H$, the phase $\epsilon_q e(\nu c_J\bar q/H)$ is a
function on $(\mathbb Z/H\mathbb Z)^\times$.  Expand it in Dirichlet
characters $\psi\pmod H$; every expansion coefficient has absolute
value at most one.  For $\nu\ne0$, consider
\[
 D_{\psi,J,\nu}(s)=\sum_{(q,H)=1}
 \frac{\lambda(q)}{q^s}\frac{\widetilde G_\nu(q)}{\sqrt q}
                 \left(\frac{\sigma JH}{q}\right)\psi(q).
\]
At $p\nmid2HJ\nu$, its local factor is
$1+\lambda(p)\xi(p)p^{-s}$, where
$\xi(q)=\psi(q)(\sigma JH\nu/q)$ is a Dirichlet character.
Multiplication by the inverse $\mathrm{GL}_2$ Euler factor leaves
\[
 1+(1-\lambda(p)^2)\xi(p)^2p^{-2s}
     +\lambda(p)\xi(p)^3p^{-3s}.
\]
These factors converge absolutely for $\Re s>1/2$.

The exceptional primes cause only a fixed power loss.  For an odd prime
$p$, put $v=v_p(\nu)$.  Direct summation over residue lifts gives
\[
 G_\nu(p^r)=
 \begin{cases}
 \varphi(p^r),&r\text{ even},\ r\le v,\\
 -p^{r-1},&r\text{ even},\ r=v+1,\\
 \epsilon_p p^{r-1/2}
       \left(\dfrac{\nu/p^{r-1}}p\right),
          &r\text{ odd},\ r=v+1,\\
 0,&\text{otherwise}.
 \end{cases}
\]
For even $r$ this is the Ramanujan sum.  For odd $r$, the sum of the
lifts vanishes unless $p^{r-1}\mid\nu$, and the primitive quadratic
Gauss sum then gives the formula.  Thus the local sum is a polynomial;
on $\Re s\ge3/4$ its absolute value is at most
$\sum_{r\ge0}(r+1)p^{-r/4}\le C_0$.
The inverse Euler polynomial is also bounded absolutely.  Primes dividing
$J$ remove the local sum altogether.  Passing to the primitive character
underlying $\xi$ introduces only such inverse Euler polynomials.  Hence
\[
 D_{\psi,J,\nu}(s)=L(s,f\otimes\xi^{\rm prim})C_{\psi,J,\nu}(s),
\]
where $C_{\psi,J,\nu}$ is holomorphic for $\Re s>1/2$ and bounded
on $\Re s\ge3/4$ by $(2J|\nu|)^C$.  Indeed the exceptional-prime
cost is at most $C_1^{\omega(J\nu)}$, and the other corrections have
summable $O(p^{-3/2})$ errors there.
The conductor of $\xi^{\rm prim}$ is at most $4HJ|\nu|$.
The twisted cusp-form $L$-function is entire; its functional equation
and polynomial strip bounds therefore give
\begin{equation}\label{odd:gauss-strip}
 |D_{\psi,J,\nu}(3/4+i t)|
 \ll_{f,H}(J|\nu|)^C(1+|t|)^C.
\end{equation}
The exponents are absolute because the degree and weight are fixed.

It remains to sum these bounds without losing the power saving.
The kernel and integration by parts in $x$ give, for fixed $r,s$ and
every $K>0$,
\begin{equation}\label{odd:weight-decay}
 |(z\partial_z)^r(R\partial_R)^s W(z,R)|
 \ll_{r,s,K}(1+z)^{-K}(1+|R|)^{-K}(1+|\log z|).
\end{equation}
The logarithm accounts for the double pole when $z$ tends to zero.
On the actual sums $z\ge X^{-1}$.
For a fixed small $\rho>0$, insert smooth cutoffs
$z\ll X^\rho$ and $|R|\ll X^\rho$.
The discarded part is negligible: use $|G_\nu(q)|\le q$ and
\[
 \sum_{\nu\in\mathbb Z}
 (1+|\nu|X/(HJq))^{-K}\ll_H1+Jq/X
\]
to obtain a polynomial absolute majorant in $X,J$ for
Equation~\eqref{odd:poisson}.  More explicitly, its absolute value is
$O_\epsilon(X^{3/2+\epsilon}\log X)$, uniformly for $J\ge1$:
the $q$-sums contribute
$(X/m_1)^{1/2+\epsilon}+(J/X)(X/m_1)^{3/2+\epsilon}$,
and the remaining $m_1\mid H^\infty$ sums converge.
For $z>1$, absorb the positive logarithm into an arbitrarily small
power of $z$ in this calculation.  On either discarded tail, increasing
$K$ in Equation~\eqref{odd:weight-decay} supplies an arbitrarily large
factor $X^{-\rho K}$.
The retained nonzero-frequency range satisfies
\begin{equation}\label{odd:ranges}
 m_1q\ll X^{1+\rho},\qquad
 0<|\nu|\ll_H JX^{2\rho},\qquad
 q\gg_H X^{1-\rho}/J.
\end{equation}

Partition the $q$-sum into smooth segments $q\asymp T$.
By Equation~\eqref{odd:weight-decay}, the finitely many derivative
norms needed for Mellin inversion of the segment's weight are
$O(X^{C'\rho})$, uniformly in $m_1,J,\nu,T$.
Mellin inversion uses $D_{\psi,J,\nu}(1+s)T^s$.
Shift its contour to $\Re s=-1/4$; the factorization above shows that
no pole is crossed.  Equation~\eqref{odd:gauss-strip} and the vertical
decay of the weight's Mellin transform give
\begin{equation}\label{odd:dyadic}
 \sum_q\frac{\lambda(q)}q\frac{\widetilde G_\nu(q)}{\sqrt q}
 \left(\frac{\sigma JH}{q}\right)\psi(q)W_*(q/T)
 \ll (J|\nu|)^C X^{C'\rho}T^{-1/4}.
\end{equation}
Here $W_*$ includes the dyadic and cutoff factors and both arguments
of the original weight; their logarithmic derivatives are controlled
by Equation~\eqref{odd:weight-decay}.
The sum over $m_1$ is bounded since
\[
 \sum_{m_1\mid H^\infty}\frac{|\lambda(m_1)|}{\sqrt{m_1}}
 \le\prod_{p\mid H}(1-p^{-1/2})^{-2}<\infty.
\]
Equation~\eqref{odd:ranges} bounds the dyadic sum of $T^{-1/4}$ by
$O_H(J^{1/4}X^{-1/4+\rho/4})$, even if its formal lower endpoint
is below one.  Also
\[
 \sum_{0<|\nu|\ll_H JX^{2\rho}}(J|\nu|)^C
 \ll_H J^{2C+1}X^{2(C+1)\rho}.
\]
Multiplying these estimates by $X/J$ from Equation~\eqref{odd:poisson}
proves $T_{\ne0}(J)\ll J^B X^{3/4+C''\rho}$.
Choose $\rho$ absolutely small enough that $C''\rho\le1/8$.
This proves Equation~\eqref{odd:offdiagonal}.

\paragraph{Conclusion.}
Sum the zero- and nonzero-frequency errors over $d\le Y$, with $J=d^2$.
Together with the large-square and residue tails, the total error is
\[
 O\left(Y^{2B+1}X^{1-\kappa}+X^{1+\rho}/Y\right),
 \qquad \kappa=1/16,
\]
after enlarging $B$.  Choose an absolute $\delta_0>0$ with
$(2B+1)\delta_0<\kappa/2$, put $Y=X^{\delta_0}$, and choose the
tail exponent $0<\rho<\delta_0/4$.
Both errors are $O(X^{1-\eta})$ for an absolute $\eta>0$.
Equation~\eqref{odd:residue} proves Equation~\eqref{odd:moment}.
Its leading coefficient is positive, so some derivative is nonzero in
every sufficiently large dyadic interval.  The odd functional equation
already forces the central value to vanish, proving the final assertion.
\end{proof}

\begin{proof}[Proof of Proposition~\ref{prop:rough-nonvanishing}\textup{(ii)}]
Choose $H$ divisible by $8N$ and by all primes in $S\cup B$.
The prescribed squareclasses and sign determine a reduced class
$a\pmod H$ with $\sigma a\equiv1\pmod4$, by the Chinese remainder
theorem. At a prime in $B\setminus S$, choose any nonzero residue.
Then every positive squarefree $n\equiv a\pmod H$ gives an odd
fundamental discriminant $D=\sigma n$ in the prescribed filter and
prime to every member of $B$. Lemma~\ref{nv:odd-first-moment} supplies
arbitrarily large such $n$ with $L'(E^{(D)},1)\ne0$.
\end{proof}
